% Generated draft: requires mathematical and source review.
\documentclass[11pt]{article}
\usepackage[margin=25mm]{geometry}
\usepackage{fontspec}
\setmainfont{Latin Modern Roman}
\usepackage{amsmath,amssymb,mathtools,unicode-math}
\setmathfont{Latin Modern Math}
\usepackage{enumitem,needspace,xurl,hyperref}
\setlength{\parindent}{0pt}
\setlength{\parskip}{5pt}
\setlength{\emergencystretch}{3em}
\newcommand{\N}{\mathbb{N}}
\newcommand{\Z}{\mathbb{Z}}
\newcommand{\Q}{\mathbb{Q}}
\newcommand{\R}{\mathbb{R}}
\newcommand{\C}{\mathbb{C}}
\newcommand{\F}{\mathbb{F}}
\newcommand{\A}{\mathbb{A}}
\newcommand{\PP}{\mathbb{P}}
\newcommand{\E}{\mathbb{E}}
\newcommand{\eps}{\varepsilon}
\newcommand{\dd}{\,\mathrm{d}}
\begin{document}
\section*{UnsolvedMath ID: 800861}
\textbf{Source problem code:} catalogue entry
\subsection{Definitions and mathematical statement}
The supplied record presents the following as Conjecture 1, concerning the type-$2$ constant of tensors. For $p\geq 2$, it asserts
\[
\mathbb{E}\left\|\sum_{i=1}^{n}g_iT_i\right\|_{\mathcal{I}_p}
\leq
\widetilde{\mathcal{O}}_{r,p}\left(
 d^{\frac12-\frac1p}
 \left(\sum_{i=1}^{n}\|T_i\|_{\mathcal{I}_p}^{2}\right)^{1/2}
\right).
\]
Here the supplied meaning of $\widetilde{\mathcal{O}}_{r,p}$ is an upper bound allowing multiplicative constants depending on $r$ and $p$, together with polylogarithmic factors in $d$ and $n$.

The supplied excerpt does not define the tensor space, the scalar field, the quantity $\|\cdot\|_{\mathcal{I}_p}$, or the roles and admissible values of $r$ and $d$. The index $n$ counts the displayed summands, but the admissible families $(T_i)_{i=1}^{n}$ are not specified. Neither the joint probability law of $(g_i)_{i=1}^{n}$ nor the probability space underlying $\mathbb{E}$ is given. In particular, Gaussianity, independence, normalization, and whether the tensors are deterministic cannot be assumed from this excerpt. The admissible range of $p$ beyond the written condition $p\geq 2$, including whether an endpoint $p=\infty$ is intended, is also unspecified. Thus the displayed assertion preserves the source formula, but its full quantified formulation remains unresolved.
\subsection{Short English statement}
The supplied conjecture bounds the expected $\mathcal{I}_p$ size of a weighted tensor sum by $d^{1/2-1/p}$ times the square root of the sum of squared individual sizes, allowing constants depending on $r,p$ and polylogarithmic factors in $d,n$. Definitions of the tensors, their size, and the random coefficients are missing, so a complete mathematical statement requires review.
\subsection{Sources}
\begin{itemize}
\item Local UnsolvedMath source registry, permanent ID $800861$, source problem code \texttt{catalogue entry}. The record attributes the problem to Bandeira et al., gives proposed year $2026$, and labels its status as open. This status is supplied registry data, not an independently verified mathematical assessment.
\item Bandeira et al., \emph{Randomstrasse101: Open Problems of 2025} ($2026$), Conjecture 1, context: type-$2$ constant of tensors. Supplied source URL: \url{https://arxiv.org/abs/2603.29571}. The record reports extraction by \texttt{arxiv\_latex} and an access date of \texttt{2026-10-16}; these are supplied provenance details, not a claim of independent source access or verification.
\end{itemize}
\end{document}
